% TEMPLATE-MILC-v3.tex - plantilla maestra UNICA de las guias MILC v3.
%
% La usan las tres familias de guias, cada una con su driver:
%   · Grados 6-11  -> scripts/build-guias-g11.py
%   · Bebras       -> scripts/build-guias-bebras.py
%   · Semillero    -> scripts/build-guias-semillero.py
%
% Antes habia tres copias casi identicas (~400 lineas iguales, 6 distintas):
% cada mejora habia que aplicarla tres veces, y en la practica no se hacia
% (el motor de recursos vivia solo en dos de ellas). Lo que varia entre
% familias esta parametrizado en 6 placeholders que cada driver llena
% (se nombran aqui SIN su sintaxis real: el builder reemplaza por texto plano,
% tambien dentro de los comentarios, y un valor multilinea rompe el preambulo):
%
%   HEADER_IZQ        encabezado izquierdo de pagina
%   HEADER_DER        encabezado derecho de pagina
%   PORTADA_ETIQUETA  rotulo sobre el numero grande (GRADO/MOMENTO/MODULO)
%   PORTADA_SERIE     linea de serie de la portada
%   PORTADA_ID        identificador de la guia en la portada
%   PORTADA_CIRCULO   el circulito de grado (°), o vacio si no aplica
%
% El resto de claves las documenta content/guias/_SCHEMA.md. El script valida
% que no queden placeholders sin reemplazar antes de escribir el archivo final.

\documentclass[11pt,letterpaper]{article}
\usepackage[margin=1.75cm]{geometry}
\usepackage[table]{xcolor}
\usepackage{fontspec}
\usepackage[spanish]{babel}
\usepackage{tikz}
% Librerías TikZ para los diagramas de las guías (bloque `recursos:`):
%  · babel  → babel en español vuelve activos caracteres como «>», y TikZ los usa
%             en sus opciones (>=stealth, ->). Sin esta librería, cualquier
%             diagrama con flechas revienta con «\language@active@arg> has an
%             extra }». Debe ir ANTES de que se dibuje nada.
%  · positioning        → sintaxis «below left=2pt and 3pt».
%  · decorations.pathreplacing → llaves ({brace}) para señalar rangos.
\usetikzlibrary{babel,positioning,decorations.pathreplacing}
\usepackage{graphicx}
% Circuitos: el trabajo de electrónica se hace hoy con micro:bit y sus sensores
% integrados (MakeCode, sin cableado), así que no se necesitan esquemáticos.
% Si algún día entran montajes con protoboard, basta descomentar esta línea:
% los diagramas .tex/.tikz declarados en `recursos:` ya se incrustan solos.
% \usepackage{circuitikz}
\usepackage[most]{tcolorbox}
\usepackage{tabularx}
\usepackage{array}
\usepackage{fancyhdr}
\usepackage{titlesec}
\usepackage{ragged2e}
\usepackage{enumitem}
\usepackage{microtype}

% Helvetica Neue no trae la flecha U+2192, asi que cada «→» escrito literalmente
% en el YAML se compilaba a NADA, en silencio (462 flechas en 61 guias: rutas del
% tipo «paleta → tipografia → lockup» salian rotas). Mapeamos el caracter a su
% version matematica, que si tiene glifo. Asi el autor puede seguir escribiendo
% «→» comodamente en el YAML.
\usepackage{newunicodechar}
\newunicodechar{→}{\ensuremath{\rightarrow}}

\setmainfont{Helvetica Neue}
\setsansfont{Helvetica Neue}
\setlength{\parindent}{0pt}
\setlength{\parskip}{6pt}
\hyphenpenalty=200
\exhyphenpenalty=200
\pretolerance=400
\tolerance=2000
\setlength{\emergencystretch}{1em}
\renewcommand{\arraystretch}{1.25}
\setlist{leftmargin=5mm,itemsep=1mm,topsep=1mm}
\newcolumntype{Y}{>{\RaggedRight\arraybackslash}X}

\definecolor{milcMagenta}{HTML}{E6007E}
\definecolor{milcVino}{HTML}{5A0038}
\definecolor{milcTurquesa}{HTML}{0093A5}
\definecolor{milcVerde}{HTML}{58A923}
\definecolor{milcMostaza}{HTML}{E5B400}
\definecolor{milcOcre}{HTML}{B86600}
\definecolor{milcNegro}{HTML}{241816}
\definecolor{milcGris}{HTML}{F7F7F4}
\definecolor{milcLinea}{HTML}{E8E2DD}
\definecolor{milcGradePrimary}{HTML}{84CC16}
\definecolor{milcGradeDark}{HTML}{4D7C0F}
\definecolor{milcGradeSoft}{HTML}{ECFCCB}
\definecolor{milcGradeAccent}{HTML}{CFE7FF}
\definecolor{milcGradeText}{HTML}{FFFFFF}
\color{milcNegro}

\pagestyle{fancy}
\fancyhf{}
\lhead{\small Tecnología e Informática · Grado 7}
\rhead{\small Guía 24 · MILC v3}
\cfoot{\small\thepage}
\renewcommand{\headrulewidth}{0pt}

\newtcolorbox{titlebox}[1]{
  enhanced,colback=#1,colframe=#1,arc=7mm,boxrule=0pt,
  left=7mm,right=7mm,top=3mm,bottom=3mm,
  before skip=8pt,after skip=8pt,
  fontupper=\sffamily\bfseries\Large\color{white}
}

\newtcolorbox{softbox}[3]{
  enhanced,colback=#2,colframe=#1,borderline west={4pt}{0pt}{#1},
  arc=2mm,boxrule=.4pt,left=4mm,right=4mm,top=3mm,bottom=3mm,
  before skip=6pt,after skip=8pt,title={#3},coltitle=white,
  colbacktitle=#1,fonttitle=\sffamily\bfseries,fontupper=\RaggedRight,
  attach boxed title to top left={xshift=3mm,yshift=-2mm},
  boxed title style={arc=2mm, boxrule=0pt}
}

\newtcolorbox{drawbox}[2]{
  enhanced,colback=white,colframe=#1,arc=4mm,boxrule=1pt,
  left=5mm,right=5mm,top=4mm,bottom=4mm,before skip=8pt,after skip=8pt,
  title={#2},coltitle=white,colbacktitle=#1,fonttitle=\sffamily\bfseries,
  fontupper=\RaggedRight,
  attach boxed title to top left={xshift=4mm,yshift=-2mm},
  boxed title style={arc=3mm, boxrule=0pt}
}

\newtcolorbox{infoband}[1]{
  enhanced,colback=#1!8,colframe=#1!50,arc=2mm,boxrule=.5pt,
  left=4mm,right=4mm,top=2mm,bottom=2mm,before skip=4pt,after skip=4pt,
  fontupper=\small\RaggedRight
}

% Puente narrativo entre secciones (contrato editorial Conéctate).
% Da continuidad: "Acabas de X. Ahora vas a Y porque Z."
% Visualmente: banda izquierda en color del grado, texto en itálica pequeña.
\newtcolorbox{puentebox}{
  enhanced,colback=milcGradeSoft,colframe=milcGradeDark,
  borderline west={3pt}{0pt}{milcGradeDark},
  arc=0mm,boxrule=0pt,left=5mm,right=4mm,top=2.5mm,bottom=2.5mm,
  before skip=4pt,after skip=8pt,
  fontupper=\itshape\small\color{milcNegro}\RaggedRight
}
% La flecha va en modo matematico a proposito: Helvetica Neue NO tiene el glifo
% U+2192, asi que \textrightarrow se compilaba a NADA («Missing character: There
% is no → in font Helvetica Neue») y el puente salia sin flecha. En modo
% matematico la flecha viene de la fuente matematica y si se dibuja.
\newcommand{\puente}[1]{%
  \begin{puentebox}%
    \textcolor{milcGradeDark}{$\rightarrow$}\hspace{2mm}#1%
  \end{puentebox}%
}

\newcommand{\linead}[1]{\noindent\color{milcLinea}\rule{#1}{0.5pt}\color{milcNegro}}
\newcommand{\respuesta}[1]{\par\vspace{2mm}\foreach \n in {1,...,#1}{\noindent\color{milcLinea}\rule{\linewidth}{0.5pt}\par\vspace{5mm}}\color{milcNegro}}
\newcommand{\checkbox}{\textcolor{milcGradeDark}{\fbox{\phantom{x}}}\hspace{5pt}}

\newcommand{\phasecard}[4]{
\begin{tcolorbox}[enhanced,colback=#3,colframe=#2,arc=3mm,boxrule=.5pt,height=3.05cm,valign=center,halign=center,left=1.5mm,right=1.5mm,top=2mm,bottom=2mm]
{\sffamily\bfseries\fontsize{22}{24}\selectfont\textcolor{#2}{#1}}\par
{\sffamily\bfseries\small #4}
\end{tcolorbox}
}

% ── Recursos visuales de la guía ──────────────────────────────────────
% Declarados en el bloque `recursos:` del YAML y emitidos por el builder.
% \guiaFigura   → imagen rasterizada o PDF (foto, carta celeste, montaje).
% \guiaDiagrama → fuente TikZ/circuitikz (.tex) incrustada como vector:
%                 es la vía para circuitos y esquemáticos editables.
% #1 = ruta relativa al .tex · #2 = pie (puede ir vacío)
\newcommand{\guiaFigura}[2]{%
\begin{center}
\includegraphics[width=0.88\linewidth,height=0.38\textheight,keepaspectratio]{#1}\\[1.5mm]
{\footnotesize\itshape\textcolor{milcNegro!75}{#2}}
\end{center}
}

\newcommand{\guiaDiagrama}[2]{%
\begin{center}
\input{#1}\\[1.5mm]
{\footnotesize\itshape\textcolor{milcNegro!75}{#2}}
\end{center}
}

\begin{document}
\thispagestyle{empty}
\begin{tikzpicture}[remember picture,overlay]
  \fill[white] (current page.south west) rectangle (current page.north east);
  \path[fill=milcGradePrimary,rounded corners=16mm]
    ([xshift=.55cm,yshift=-.55cm]current page.north west) rectangle
    ([xshift=-.55cm,yshift=3.65cm]current page.south east);

  \node[anchor=north west,text=milcGradeText] at ([xshift=2.0cm,yshift=-1.85cm]current page.north west)
    {\sffamily\bfseries\fontsize{14}{16}\selectfont GRADO};
  \node[anchor=north west,text=milcGradeText,opacity=.82] at ([xshift=2.0cm,yshift=-2.75cm]current page.north west)
    {\sffamily\bfseries\fontsize{9}{11}\selectfont SERIE GUÍAS · TECNOLOGÍA E INFORMÁTICA};

  \begin{scope}[shift={([xshift=-3.05cm,yshift=-2.55cm]current page.north east)}]
    \fill[white,opacity=.80] (-.62,-.52) rectangle (.62,.52);
    \draw[milcGradeText,line width=1pt] (-.62,-.52) rectangle (.62,.52);
    \fill[milcGradeDark] (-.42,-.38) rectangle (-.22,.06);
    \fill[milcGradeAccent] (-.10,-.38) rectangle (.10,.24);
    \fill[milcGradePrimary!60!black] (.22,-.38) rectangle (.42,.38);
  \end{scope}

  \node[anchor=west,text=milcGradeText] at ([xshift=1.9cm,yshift=-12.85cm]current page.north west)
    {\sffamily\bfseries\fontsize{124}{124}\selectfont 7};
  % El circulito de grado (°) solo aplica a las guias de grado («11°»); en
  % Bebras y Semillero el numero grande es un momento/modulo, no un grado.
  \draw[milcGradeText,line width=1.7pt]
    ([xshift=4.52cm,yshift=-11.68cm]current page.north west) circle (.13cm);

  \node[anchor=west,text=milcGradeText,text width=8.7cm,align=left] at ([xshift=10.35cm,yshift=-11.6cm]current page.north west)
    {
     {\sffamily\bfseries\fontsize{19}{22}\selectfont Cómo aprende\\una IA ---\\datos,\\entrenamiento\\y sesgo}\\[4mm]
     {\sffamily\bfseries\fontsize{23}{26}\selectfont Guía 24-7-TIC}\\[2mm]
     {\sffamily\fontsize{13}{16}\selectfont Período 3 · Inteligencia Artificial}\\[5mm]
     {\sffamily\bfseries\fontsize{11}{13}\selectfont Institución Educativa Sor María Juliana}\\[1mm]
     {\sffamily\fontsize{10}{12}\selectfont Autor: Dr. Álvaro Cárdenas Orozco}
    };

  \path[fill=milcGradeSoft,rounded corners=7mm]
    ([xshift=1.35cm,yshift=.85cm]current page.south west) rectangle
    ([xshift=-1.35cm,yshift=4.25cm]current page.south east);
  \draw[milcGradeDark,line width=.65pt,rounded corners=7mm]
    ([xshift=1.35cm,yshift=.85cm]current page.south west) rectangle
    ([xshift=-1.35cm,yshift=4.25cm]current page.south east);
  \node[anchor=center,text=milcNegro,text width=17.0cm,align=center] at ([yshift=2.55cm]current page.south)
    {\sffamily\fontsize{10}{12}\selectfont Metodología MILC · Apertura ancestral · Pensamiento computacional · Triángulo Dussel-Estoicismo-Floridi\\
    \textcolor{milcGradeDark}{\bfseries Producto:} En el cuaderno: \textbf{(a) diagrama del aprendizaje de una IA} con 4 partes (datos → algoritmo → modelo → predicción); \textbf{(b) 3 ejemplos reales de sesgo en IA} con análisis (qué pasó, por qué, consecuencias); \textbf{(c) tu propio mini-ejemplo} de cómo un sesgo en datos puede causar problema (puedes inventar un caso simple); \textbf{(d) 5 reglas para detectar sesgo en una IA que uses}. Firma.};
\end{tikzpicture}
\mbox{}
\newpage

\section*{Apertura: Cómo aprende una IA --- datos, entrenamiento y el problema del sesgo}

\begin{softbox}{milcGradeDark}{milcGradeSoft}{Saber ancestral}
\textbf{Cuando doña Mercedes, maestra rural del Valle del Cauca, enseñaba a leer a un niño nuevo, no le mostraba 1 sola palabra: le mostraba cientos.} Le mostraba \emph{``casa''} y la repetía. Le mostraba \emph{``mamá''}, \emph{``perro''}, \emph{``árbol''}. Cada palabra varias veces. Con el tiempo, el niño \textbf{empezaba a reconocer patrones}: \emph{las letras juntas forman palabras, las palabras tienen sonidos, los sonidos forman frases}. \emph{Doña Mercedes no le ``programaba'' las reglas}: el niño las extraía de los ejemplos. \textbf{Esa es exactamente la forma en que aprende una IA moderna}: le muestras miles de ejemplos, ella encuentra patrones, y después puede generalizar. \emph{Si los ejemplos son buenos y variados, aprende bien. Si los ejemplos son malos o sesgados, aprende mal}. Como doña Mercedes: si solo le hubiera mostrado al niño palabras de un grupo (solo objetos urbanos, sin animales del campo), el niño habría aprendido un mundo parcial. Las IAs igual.
\end{softbox}

\begin{softbox}{milcTurquesa}{milcGris}{Saber contemporáneo}
Una IA \emph{aprende} no porque le programen reglas, sino porque \textbf{le muestran muchos ejemplos}. Este proceso se llama \textbf{entrenamiento} y tiene 4 partes: \textbf{(1) Datos de entrenamiento:} miles o millones de ejemplos. Para ChatGPT, son libros, artículos, sitios web, foros. Para visión por computador, son fotos etiquetadas. Para reconocimiento facial, son rostros. \textbf{(2) Algoritmo de aprendizaje:} la fórmula matemática que ajusta el modelo según los datos. El más popular: \emph{redes neuronales profundas} (deep learning). \textbf{(3) Modelo entrenado:} el resultado del entrenamiento. \emph{Una IA en sí misma}: un conjunto de números que codifican lo aprendido. \textbf{(4) Predicción/generación:} cuando le das algo nuevo, el modelo lo procesa y genera respuesta. \textbf{El problema del sesgo (bias)}: si los datos de entrenamiento están \emph{sesgados} (incompletos, parciales, con prejuicios), la IA aprende esos sesgos y los reproduce. Por ejemplo, si una IA de RH solo vio currículos de hombres, va a discriminar mujeres. Esto es real y ha pasado en empresas grandes.
\end{softbox}

\begin{softbox}{milcMostaza}{milcGris}{Pregunta-puente}
\textit{Si una IA tiene que aprender a reconocer ``buenas notas'', pero la entrenan SOLO con notas de niños de colegios privados de Bogotá, ¿\textbf{qué pasa cuando le muestran una nota} de un niño del campo del Pacífico que escribe con caligrafía distinta?}
\end{softbox}

\begin{softbox}{milcVerde}{milcGris}{Saber y saber hacer}
Al terminar la clase: (1) podrás \textbf{identificar} las 4 partes del aprendizaje de una IA; (2) sabrás \textbf{explicar} qué es el sesgo y por qué ocurre; (3) podrás \textbf{analizar} 3 ejemplos reales de sesgo; (4) habrás \textbf{creado} tu diagrama del aprendizaje + 5 reglas para detectar sesgo.
\end{softbox}



\small
\begin{tabularx}{\textwidth}{>{\bfseries}p{3.0cm}Y}
\rowcolor{milcGradePrimary!12}Autor & Dr. Álvaro Cárdenas Orozco\\
DBA & Reconoce los fundamentos, usos y límites éticos de la Inteligencia Artificial (MEN, T\&I 7°)\\
Referentes & Saber ancestral del consejero · Modelos de lenguaje · Ética algorítmica y prompting\\
Producto final & En el cuaderno: \textbf{(a) diagrama del aprendizaje de una IA} con 4 partes (datos → algoritmo → modelo → predicción); \textbf{(b) 3 ejemplos reales de sesgo en IA} con análisis (qué pasó, por qué, consecuencias); \textbf{(c) tu propio mini-ejemplo} de cómo un sesgo en datos puede causar problema (puedes inventar un caso simple); \textbf{(d) 5 reglas para detectar sesgo en una IA que uses}. Firma.\\
\end{tabularx}
\normalsize

\puente{Hoy haces 4 cosas: descubres cómo aprende una IA, analizas 3 casos de sesgo, dibujas el diagrama, formulas las 5 reglas.}

\begin{titlebox}{milcGradeDark}
Ruta MILC: así vamos a aprender
\end{titlebox}
\begin{tabularx}{\textwidth}{>{\centering\arraybackslash}X>{\centering\arraybackslash}X>{\centering\arraybackslash}X>{\centering\arraybackslash}X}
\phasecard{1}{milcVerde}{milcGris}{Escucha\\[-1mm]\small Observo, escucho y pregunto.} &
\phasecard{2}{milcTurquesa}{milcGris}{Sistematización\\[-1mm]\small Pensamiento computacional.} &
\phasecard{3}{milcMagenta}{milcGris}{Praxis\\[-1mm]\small Construyo y aplico.} &
\phasecard{4}{milcVino}{milcGris}{Evaluación\\[-1mm]\small Reflexión liberadora.}
\end{tabularx}


\puente{Empezamos con un experimento mental: enseñar a una IA a reconocer una palabra.}

\begin{titlebox}{milcVerde}
Fase 1 · Escucha
\end{titlebox}

\begin{softbox}{milcVerde}{milcGris}{Escena para escuchar}
\textbf{Actividad 1 · ANALIZA --- Enseñar a una IA a reconocer ``gato''} (10 min · individual). Imagina este experimento: \emph{quieres entrenar una IA para que reconozca fotos de gatos. Le muestras 1000 fotos de gatos para que aprenda}. En el cuaderno responde: \emph{(1) ¿Bastan 1000 fotos? ¿O necesitas más?} \emph{(2) ¿Qué tipo de fotos deberías incluir? (todas iguales o variadas)}. \emph{(3) ¿Qué pasa si todas las fotos son de gatos negros? ¿La IA reconocerá gatos blancos?}. \emph{(4) ¿Qué pasa si solo le muestras fotos de gato y NO de otras cosas? ¿La IA podrá distinguir gato de perro?}. \emph{(5) ¿Quién decide qué fotos usar para entrenar?}.
\end{softbox}

\checkbox Imaginé el experimento de entrenar la IA.\\
\checkbox Respondí las 5 preguntas con calma.\\
\checkbox Identifiqué que los datos importan, no solo cantidad.

\begin{infoband}{milcVerde}


\textbf{Lo que va al cuaderno (Actividad 1 · ANALIZA).} Título: «Actividad 1 --- Entrenar una IA: la importancia de los datos». \textbf{Formato:} 5 respuestas cortas. \textbf{Extensión:} 5-6 líneas. \textbf{Sabes que terminaste cuando} entiendes que los datos buenos son tan importantes como el algoritmo.
\end{infoband}

\puente{Después aprendes las 4 partes del proceso + 3 casos reales de sesgo.}

\begin{titlebox}{milcTurquesa}
Fase 2 · Sistematización con pensamiento computacional
\end{titlebox}

\begin{softbox}{milcTurquesa}{milcGris}{Cuatro pilares aplicados al tema}
Tu intuición es correcta: \textbf{los datos son lo más importante en el aprendizaje de una IA}. Ahora aprende las 4 partes del proceso + 3 casos reales de sesgo.
\end{softbox}

\small
\begin{tabularx}{\textwidth}{>{\bfseries\RaggedRight\arraybackslash}p{3.6cm}Y}
\rowcolor{milcTurquesa!90}\color{white}Pilar & \color{white}\bfseries Aplicación al tema\\
1. Descomposición & \textbf{Las 4 partes del aprendizaje de una IA.} \textbf{(1) Datos de entrenamiento.} Miles o millones de ejemplos. Para ChatGPT: textos de libros, Wikipedia, foros, sitios web. Para una IA médica: historiales clínicos. Para reconocimiento facial: fotos de personas con sus nombres. \textbf{Calidad de los datos = calidad de la IA}. Datos sesgados → IA sesgada. \textbf{(2) Algoritmo de aprendizaje.} La fórmula matemática que ajusta el modelo. La más usada hoy: \emph{redes neuronales profundas} (deep learning), inspiradas en cómo funcionan las neuronas del cerebro humano. Otras: árboles de decisión, regresión, máquinas de soporte vectorial.\\
2. Reconocimiento de patrones & \textbf{Las 4 partes (continuación).} \textbf{(3) Modelo entrenado.} El resultado del proceso: un conjunto de \emph{parámetros} (números) que el algoritmo ajustó. \emph{La IA misma}. GPT-4 tiene aproximadamente 1.7 trillones de parámetros. Cada uno es un número ajustado durante el entrenamiento. \textbf{(4) Predicción/generación.} Cuando le das una entrada nueva, el modelo la procesa y produce una salida. \emph{El uso real}. Para ChatGPT: tú escribes pregunta, modelo genera respuesta. Para reconocedor facial: tú muestras foto, modelo dice quién es.\\
3. Abstracción & \textbf{Qué es el sesgo (bias) y por qué ocurre.} El \textbf{sesgo} en IA es cuando la IA discrimina o falla sistemáticamente con cierto grupo. \emph{No es un bug accidental: es consecuencia de los datos}. Si entrenaste la IA con datos parciales, la IA aprenderá esa parcialidad. Es como un niño criado solo en la ciudad: si lo llevas al campo, no sabe nombrar animales. \textbf{3 ejemplos reales de sesgo:} (a) \textbf{Amazon (2018):} desarrolló IA para filtrar CV y la entrenó con CVs de 10 años anteriores (mayoritariamente hombres). La IA descartaba CVs de mujeres automáticamente. Amazon la canceló. (b) \textbf{Google Photos (2015):} clasificó fotos de personas afrodescendientes como \emph{``gorillas''}. Los datos no incluían suficiente diversidad. Google se disculpó públicamente. (c) \textbf{COMPAS (sistema de justicia EEUU):} IA usada para evaluar riesgo de reincidencia mostró sesgo contra personas afroamericanas. Influyó en condenas reales. Sigue siendo controversial.\\
4. Algoritmo conceptual & \textbf{Las 5 reglas para detectar sesgo en una IA.} (1) \textbf{Pregunta de qué datos aprendió.} Si no te lo dicen, sospecha. (2) \textbf{Verifica resultados en distintos grupos.} ¿La IA funciona igual de bien con hombres y mujeres? ¿Con personas de distinta edad/origen? Si trata distinto a grupos similares, hay sesgo. (3) \textbf{Busca tu propio ejemplo de prueba.} Si vives en zona rural, prueba la IA con ejemplos rurales y ve cómo responde. (4) \textbf{Lee la documentación.} Las IAs profesionales explican (o deberían explicar) sus limitaciones y sesgos conocidos. (5) \textbf{No asumas que la IA es ``neutral''.} \emph{Toda IA refleja los valores y limitaciones de sus creadores y datos}. La neutralidad pura no existe.\\
\end{tabularx}
\normalsize

\begin{drawbox}{milcTurquesa}{4 partes del aprendizaje + sesgo}
\textbf{4 partes:} \\[1mm]
\textbf{1.} DATOS (miles de ejemplos). \\[1mm]
\textbf{2.} ALGORITMO (fórmula que ajusta). \\[1mm]
\textbf{3.} MODELO (resultado entrenado). \\[1mm]
\textbf{4.} PREDICCIÓN/GENERACIÓN (uso real). \\[1mm]
\textbf{Sesgo:} cuando datos parciales generan IA discriminatoria. Real, no teórico.
\end{drawbox}

\begin{softbox}{milcMostaza}{milcGris}{Errores comunes y su corrección}
\textbf{Mitos sobre cómo aprende la IA.} (1) \emph{``La IA es objetiva y no tiene prejuicios''} --- Falso. La IA refleja prejuicios de sus datos. (2) \emph{``Si la IA es nueva, no tiene sesgo''} --- Falso. Las más nuevas también tienen sesgo, solo que distinto. (3) \emph{``Más datos = sin sesgo''} --- Falso. Si los datos son sesgados, más datos también pueden estar sesgados. (4) \emph{``Solo los datos viejos están sesgados''} --- Falso. Datos nuevos también pueden tener sesgos contemporáneos. (5) \emph{``El sesgo es problema técnico que se arregla solo''} --- Falso. Es problema técnico Y social que requiere atención humana constante.
\end{softbox}

\begin{infoband}{milcTurquesa}


\textbf{Lo que va al cuaderno (Actividad 2 · EXPLICA).} Título: «Actividad 2 --- 4 partes + sesgo». \textbf{Formato:} bloque A con diagrama de las 4 partes (flechas: datos → algoritmo → modelo → predicción). Bloque B con tabla de 3 casos de sesgo (Amazon, Google Photos, COMPAS). \textbf{Extensión:} 12 líneas. \textbf{Sabes que terminaste cuando} podrías explicarle a tu hermano qué es sesgo en IA con un ejemplo real.
\end{infoband}

\puente{Con todo claro, armas tu diagrama + 3 casos + 5 reglas + mini-ejemplo propio.}

\begin{titlebox}{milcMagenta}
Fase 3 · Praxis con pensamiento computacional
\end{titlebox}

\begin{softbox}{milcMagenta}{milcGris}{Construyendo el producto con los cuatro pilares}
\textbf{Actividad 3 · CREA --- Diagrama + análisis + mini-ejemplo + 5 reglas} (32 min · individual). Vas a sintetizar todo lo aprendido en un solo trabajo.
\end{softbox}

\small
\begin{tabularx}{\textwidth}{>{\bfseries\RaggedRight\arraybackslash}p{3.6cm}Y}
\rowcolor{milcMagenta!90}\color{white}Pilar & \color{white}\bfseries Aplicación al producto\\
Descomposición & \textbf{Paso 1 --- Dibuja el diagrama del aprendizaje.} En una hoja del cuaderno, dibuja un diagrama horizontal con \textbf{4 cajas conectadas por flechas}: \emph{Datos → Algoritmo → Modelo → Predicción}. Adentro de cada caja, una descripción breve (1-2 líneas) de qué pasa ahí.\\
Patrones & \textbf{Paso 2 --- Análisis de los 3 casos reales.} Debajo del diagrama, tabla de \textbf{3 filas y 4 columnas}. Columnas: \emph{(1) Caso}, \emph{(2) ¿Qué pasó?}, \emph{(3) ¿Por qué? (sesgo en datos)}, \emph{(4) ¿Cuál fue la consecuencia?}. Filas: Amazon CV, Google Photos, COMPAS.\\
Abstracción & \textbf{Paso 3 --- Tu mini-ejemplo de sesgo.} Inventa un caso simple. Ejemplo: \emph{``Una IA para clasificar canciones colombianas se entrena solo con vallenato. Después se le muestra una salsa choke. ¿Qué pasa? La IA dice 'no es música colombiana' porque no aprendió otros géneros''}. Describe tu caso en 3-4 líneas: \emph{(a) qué se quiere lograr}, \emph{(b) qué datos usaron}, \emph{(c) cómo se manifiesta el sesgo}, \emph{(d) qué solución propondrías}.\\
Algoritmo de armado & \textbf{Paso 4 --- Las 5 reglas para detectar sesgo.} Lista numerada del 1 al 5. Reproduce las 5 reglas que aprendiste, con tus palabras. Subráyalas: serán tu kit mental para evaluar cualquier IA que uses.\\
\end{tabularx}
\normalsize

\begin{drawbox}{milcMagenta}{Antes de cerrar la S4}
\checkbox Tengo diagrama de las 4 partes del aprendizaje. \\[1mm]
\checkbox La tabla de 3 casos reales está completa. \\[1mm]
\checkbox Tengo mi mini-ejemplo de sesgo inventado. \\[1mm]
\checkbox Las 5 reglas para detectar sesgo están subrayadas. \\[1mm]
\checkbox Apliqué las reglas a 1 IA real y detecté 1 posible sesgo. \\[1mm]
\checkbox Hay reflexión personal firmada.
\end{drawbox}

\begin{softbox}{milcMagenta}{milcGris}{Plantilla de trabajo}
\textbf{Estructura.} \textit{Paso 1:} diagrama 4 cajas (datos → algoritmo → modelo → predicción). \textit{Paso 2:} tabla 3 casos reales (Amazon, Google Photos, COMPAS). \textit{Paso 3:} mini-ejemplo propio inventado. \textit{Paso 4:} 5 reglas para detectar sesgo. \textit{Paso 5:} aplicar reglas a 1 IA real. \textit{Paso 6:} reflexión + firma.
\end{softbox}

\begin{infoband}{milcMagenta}


\textbf{Lo que va al cuaderno (Actividad 3 · CREA).} Título: «Actividad 3 --- Cómo aprende una IA + sesgo». \textbf{Formato:} diagrama + 3 casos + mini-ejemplo + 5 reglas + aplicación + reflexión. \textbf{Extensión:} 1 página y media. \textbf{Sabes que terminaste cuando} podrías ser ``crítico de IA'' al revisar cualquier app.
\end{infoband}

\puente{El producto es ese diagrama + análisis + 5 reglas firmadas.}

\begin{titlebox}{milcMostaza}
Producto de la sesión
\end{titlebox}
\textbf{Diagrama del aprendizaje + 3 casos de sesgo + mini-ejemplo + 5 reglas · firmado}.
\begin{softbox}{milcMostaza}{milcGris}{Criterios de calidad}
(1) El diagrama tiene las 4 partes correctamente conectadas. (2) La tabla de 3 casos reales está completa (qué pasó, por qué, consecuencia). (3) Hay mini-ejemplo de sesgo inventado por el estudiante. (4) Las 5 reglas para detectar sesgo están aplicadas a 1 IA real. (5) Hay reflexión personal sobre el descubrimiento del sesgo.
\end{softbox}

\puente{Antes de salir, verifica que tu mini-ejemplo muestra cómo el sesgo entra por los datos.}

\begin{titlebox}{milcVino}
Fase 4 · Evaluación liberadora — cinco dimensiones
\end{titlebox}

\begin{softbox}{milcVino}{milcGris}{Sentido de la evaluación}
Evaluar no es solo revisar una nota. Es reconocer qué aprendí, qué cambió en mí, qué emociones manejé, cómo me relaciono con la ciudad y la comunidad, y qué le dejaré a quien viene detrás.
\end{softbox}

\textbf{Mi reflexión liberadora en cinco dimensiones.} Completa cada frase iniciada:

\small
\begin{tabularx}{\textwidth}{>{\bfseries\RaggedRight\arraybackslash}p{3.4cm}Y>{\RaggedRight\arraybackslash}p{4.6cm}}
\rowcolor{milcVino}\color{white}Dimensión & \color{white}\bfseries Mi respuesta (frame starter) & \color{white}\bfseries Algo que descubrí\\
1. Personal & Lo que más cambió en mí fue \linead{6.5cm} & \linead{4.2cm}\\
2. Emocional & La emoción más fuerte fue \linead{2.5cm} y la regulé \linead{3cm} & \linead{4.2cm}\\
3. Ciudadana & En democracia esto importa porque \linead{6cm} & \linead{4.2cm}\\
4. Local (Cartago) & En mi barrio sería útil para \linead{6.5cm} & \linead{4.2cm}\\
5. Intergeneracional & A mi abuela/abuelo le diría \linead{6.5cm} & \linead{4.2cm}\\
\end{tabularx}
\normalsize

\begin{infoband}{milcVino}
\textbf{Para la próxima sustentación.} Vuelve a esta tabla en seis meses. Verás que las respuestas cambian — no porque lo aprendido cambie, sino porque tú cambias. La evaluación liberadora es una conversación contigo a través del tiempo.
\end{infoband}


\puente{Cerramos con 3 voces sobre por qué entender el sesgo es entender a la IA por dentro.}

\begin{titlebox}{milcGradeDark}
Triángulo de pensamiento · cierre filosófico
\end{titlebox}

\begin{softbox}{milcVino}{milcGris}{Dussel · lente del nosotros}
\textit{````La IA refleja a quien la creó. Si los datos vinieron de unos pocos, la IA reproducirá los valores de unos pocos. La diversidad de datos es justicia.''''}

\textbf{La mayoría de IAs del mundo se entrenan con datos en inglés, de Estados Unidos y Europa}. Cuando esas IAs llegan a Colombia, traen consigo \emph{las miradas culturales de sus creadores}. Pueden tener sesgos contra acentos latinos, contra apellidos hispanos, contra realidades de barrios populares. \emph{Eso no es accidente; es consecuencia de qué datos usaron y quién los recolectó}. La justicia algorítmica empieza por diversificar quién crea las IAs y con qué datos.

\textbf{¿Las IAs que uso fueron creadas pensando en mí, mi familia, mi región?}
\end{softbox}
\linead{\linewidth}\\[1mm]
\linead{\linewidth}

\begin{softbox}{milcMostaza}{milcGris}{Estoicismo · Marco Aurelio · cuidado interior}
\textit{````Saber cómo aprende la IA te quita el asombro mágico. Lo mágico cede paso a lo analítico, y eso es ganancia.''''}

\textbf{Antes de hoy, las respuestas de ChatGPT te parecían casi mágicas}: \emph{¿cómo puede saber tanto?}. Hoy entiendes: \emph{aprendió de miles de millones de textos durante meses de entrenamiento, ajustando trillones de parámetros}. \emph{Es resultado de procesamiento matemático masivo, no magia}. Esa comprensión te coloca en posición de analizar, no de venerar.

\textbf{¿Antes veía la IA como magia? ¿Cómo cambia mi relación al saber cómo aprende?}
\end{softbox}
\linead{\linewidth}\\[1mm]
\linead{\linewidth}

\begin{softbox}{milcTurquesa}{milcGris}{Floridi · lente de la infoesfera}
\textit{````El sesgo no es defecto técnico que arreglar después: es valor político que debe debatirse antes. La IA refleja decisiones humanas.''''}

\textbf{El sesgo en IA no es problema solo técnico}: es problema \emph{político}. Detrás de cada decisión sobre qué datos usar, qué algoritmo escoger, qué medir como ``buen resultado'' \emph{hay personas que decidieron}. La IA refleja esas decisiones. Por eso debatir IA es debatir poder: \emph{quién la hace, para quién, con qué datos, para qué fines}. Esa conversación apenas empieza en el mundo. Tu generación va a ser parte de ella.

\textbf{¿Qué decisiones políticas estoy aceptando sin verlas cada vez que uso una IA?}
\end{softbox}
\linead{\linewidth}\\[1mm]
\linead{\linewidth}

\begin{titlebox}{milcVino}
Compromiso final
\end{titlebox}

\small
\begin{tabularx}{\textwidth}{>{\RaggedRight\arraybackslash}p{3.0cm}YYY}
\rowcolor{milcVino}\color{white}\bfseries Criterio & \color{white}\bfseries Lo logré & \color{white}\bfseries En proceso & \color{white}\bfseries Necesito apoyo\\
Comprensión del tema & Explico ideas centrales con mis palabras. & Reconozco ideas pero falta precisión. & Necesito repasar conceptos base.\\
Pensamiento computacional & Apliqué los 4 pilares al diseñar mi producto. & Apliqué algunos pilares. & Necesito releer la fase 2.\\
Aplicación y producto & Mi producto cumple los criterios y se puede revisar. & Producto funcional con ajustes pendientes. & Aún en construcción.\\
Reflexión 5 dimensiones & Respondí cada dimensión con honestidad. & Respondí algunas con detalle. & Mis respuestas son muy generales.\\
Triángulo de pensamiento & Conecté las 3 voces con mi trabajo real. & Conecté al menos una. & No relaciono las voces con mi proceso.\\
\end{tabularx}
\normalsize

\textbf{Hoy entendí que:}\\[1mm]
\linead{\linewidth}\\[1mm]
\linead{\linewidth}

\textbf{Mi compromiso para la próxima sesión es:}\\[1mm]
\linead{\linewidth}\\[1mm]
\linead{\linewidth}

\begin{softbox}{milcMostaza}{milcGris}{Cita de despedida · Marco Aurelio}
\textit{``El obstáculo en el camino se vuelve el camino. Lo que estorba al camino se vuelve camino.''}\\[1mm]
Cada error de hoy, bien observado, se convierte en pieza del camino que pisarás mañana.
\end{softbox}

\small
\begin{tabularx}{\textwidth}{>{\bfseries}p{3.6cm}Y}
\rowcolor{milcGradeDark!12}Nombre completo: & \linead{10cm}\\
Fecha: & \linead{4cm} \quad Curso/grupo: \linead{4cm}\\
Mi firma: & \linead{9cm}\\
\end{tabularx}
\normalsize

\begin{infoband}{milcGradeDark}
\textbf{Para la próxima sesión.} Esta semana, cuando uses una IA, pregúntate: \emph{``¿de qué datos aprendió esto? ¿hay grupos a los que no representa?''}. La próxima clase aprendes a fondo los \textbf{Modelos de Lenguaje (LLMs)}: ChatGPT, Claude, Gemini. Vas a abrir uno y conversar por primera vez.
\end{infoband}

\end{document}
